% !PRINTED 462
\weberpage{462}
%<<CONT
$\alpha, \alpha_1, \alpha_2 \ldots \alpha_{n-1}$ von $f(x) = 0$ verschwinden, d.\,h. es ist
\[
\varphi(\Theta) = 0,\quad \varphi(\Theta_1) = 0,\quad \varphi(\Theta_2) = 0 \ldots \varphi(\Theta_{n-1}) = 0.
\]

Wenn also die mit $\Theta$ conjugirten Werthe alle von einander
verschieden sind, so ist $\varphi(t)$ mit $\Phi(t)$ identisch, d.\,h.
$\Phi(t)$ ist irreducibel. Sind aber unter den mit $\Theta$
conjugirten Werthen nur $n_1$ von einander verschiedene vorhanden,
etwa $\Theta, \Theta_1, \Theta_2 \ldots \Theta_{n_1-1}$, so ist
\weq{10}{\varphi(t) = (t - \Theta)(t - \Theta_1) \ldots (t - \Theta_{n_1-1}),}
und jeder andere irreducibele Factor von $\Phi(t)$ ist, da er
wenigstens für einen der conjugirten Werthe $\Theta$ und folglich
für alle verschwinden muss, mit $\varphi(t)$ identisch. Es ist also
$\Phi(t)$ eine Potenz von $\varphi(t)$, etwa
\weq{11}{\Phi(t) = \varphi(t)^{n_2}}
und
\weq{12}{n = n_1 n_2.}
Daraus ergiebt sich also der Satz:

\gesperrt{2. Die Function $\Phi(t)$ ist entweder irreducibel oder sie
ist eine Potenz einer irreducibeln Function. Die $n$ mit einer Zahl
in $\Omega(\alpha)$ conjugirten Zahlen sind entweder alle von
einander verschieden, oder sie zerfallen in $n_1$ Systeme von je
$n_2$ unter einander gleichen Zahlen. Im ersten Falle ist $\Phi(t)$
irreducibel, im zweiten die $n_2^{\mathrm{te}}$ Potenz einer
irreducibeln Function $n_1^{\mathrm{ten}}$ Grades in $\Omega$.}

Eine Grösse $\Theta$ in $\Omega(\alpha)$, die von allen ihren
conjugirten Zahlen verschieden ist, und die also einer
irreducibeln Gleichung $n^{\mathrm{ten}}$ Grades genügt, heisst eine
\gesperrt{primitive Grösse des Körpers}. Nach dem Satz \S.\,143, 1. lassen sich
unendlich viele solche primitive Grössen bestimmen, sogar so, dass
die Co\"efficienten von $\Theta = \chi(\alpha)$ rationale Zahlen sind.
Man braucht nur über die Co\"efficienten von $\chi$ so zu verfügen,
dass unter den conjugirten Grössen $\chi(\alpha)$ keine gleichen
vorkommen.

\gesperrt{3. Jede Grösse $\omega$ des Körpers $\Omega(\alpha)$ kann
rational durch eine beliebige primitive Grösse $\Theta$ des Körpers
ausgedrückt werden.}

Denn sind $\omega, \omega_1, \omega_2 \ldots \omega_{n-1}$ die zu
$\omega$ conjugirten Zahlen, ebenso $\Theta, \Theta_1, \Theta_2
\ldots \Theta_{n-1}$ die zu $\Theta$ conjugirten und
%CONT>>
